\documentclass[11pt]{scrartcl}

\usepackage[magyar]{babel}
\usepackage[stixtwo]{fontsetup}
\usepackage{amsmath}
\usepackage[margin=1in]{geometry}
\usepackage{tikz}
\usepackage{pgfplots}
\usepackage{luacode}

\usetikzlibrary{arrows.meta}
\pgfplotsset{compat=1.18}

% Feliratméretek. A weboldalon (lwarp-pal konvertálva) az ábrák kicsinyítve
% jelennek meg, ott nagyobb feliratok kellenek, mint a PDF-ben.
\makeatletter
\@ifpackageloaded{lwarp}{%
	\newcommand{\figfont}{\large}%
	\newcommand{\figtickfont}{\large}%
}{%
	\newcommand{\figfont}{\normalsize}%
	\newcommand{\figtickfont}{\small}%
}
\makeatother

\setlength{\parindent}{0pt}
\setlength{\parskip}{0.66em}

\definecolor{fcolor}{HTML}{1F77B4}
\definecolor{gcolor}{HTML}{D62728}

\pgfplotsset{
	fn/.style={
			axis lines=middle,
			axis line style={-{Stealth[scale=1.1]}, gray!70!black, line width=1pt},
			xlabel={$x$},
			ylabel={$y$},
			xlabel style={anchor=west},
			ylabel style={anchor=south},
			tick style={gray!70!black},
			ticklabel style={font=\figtickfont, fill=white, inner sep=1pt},
			label style={font=\figfont},
			% A tengelyek és a feliratok (fehér háttéren) a görbék fölött.
			axis on top,
			width=\linewidth,
			height=7.5cm,
			xmin=\wxmin, xmax=\wxmax,
			ymin=\wymin, ymax=\wymax,
			clip=false,
		},
	% Nagyított ábra: a tengelyek a bal és az alsó szélen.
	zoom/.style={
			fn,
			axis lines=left,
			width=0.8\linewidth,
		},
}
\tikzset{
	% A TikZ alapértelmezésénél vastagabb vonalak: a weboldalon az ábrák
	% kicsinyítve jelennek meg.
	every picture/.append style={line width=0.9pt, font=\figfont},
	% Ha egy görbe kifut az ábrából, ott nyílhegyben végződik; ha valóban
	% véget ér, ott pontban (end dot); a kimaradó pontok üres karikák (hole).
	tip/.tip={Stealth[length=8pt, width=7pt]},
	end dot/.style={only marks, mark=*, mark size=2.8pt},
	hole/.style={only marks, mark=*, mark size=2.8pt,
			mark options={fill=white, line width=1.4pt}},
	graph/.style={line width=2.2pt, line cap=round},
	mirror/.style={densely dashed, line width=1pt, gray!75},
}

% A grafikonokat Lua rajzolja: a függvényt darabonként mintavételezi, ahol a
% görbe kilép az ablakból, ott (felezéssel) pontosan a szélen vágja el, és
% nyílheggyel zárja; a szakadási helyeken maga teszi ki a pontokat és a
% karikákat.
\begin{luacode*}
plot = { fn = {} }

-- EPS: ilyen közelről nézzük a töréspontokat és a végpontokat. Ha túl
-- kicsi, a 0/0 alakú kifejezések (pl. cos(3x) - 1) számítása pontatlan.
local EPS

function plot.window(xmin, xmax, lo, hi)
  plot.xmin, plot.xmax, plot.lo, plot.hi = xmin, xmax, lo, hi
  EPS = 1e-6 * (xmax - xmin)
end

local function visible(y)
  return y == y and y >= plot.lo and y <= plot.hi
end

-- x0-ban a görbe látszik, x1-ben nem: hol lép ki az ablakból?
local function crossing(f, x0, x1)
  local edge = f(x1) < plot.lo and plot.lo or plot.hi
  for _ = 1, 60 do
    local m = (x0 + x1) / 2
    if visible(f(m)) then x0 = m else x1 = m end
  end
  return { x0, edge }
end

local function emit(style, points, startTip, endTip)
  if #points < 2 then return end
  local out = {}
  for _, p in ipairs(points) do
    out[#out + 1] = string.format('(%.5f,%.5f)', p[1], p[2])
  end
  tex.sprint(string.format('\\addplot[graph, %s, %s-%s] coordinates {%s};',
    style, startTip and 'tip' or '', endTip and 'tip' or '',
    table.concat(out, ' ')))
end

-- A görbe az (a, b) nyílt intervallumon, amelyen belül nincs pólus.
local function curve(f, a, b, style, leftTip, rightTip)
  local n = math.max(60, math.ceil(1500 * (b - a) / (plot.xmax - plot.xmin)))
  local run, startTip, px
  for i = 0, n do
    local x = i == 0 and a + EPS or i == n and b - EPS or a + (b - a) * i / n
    local y = f(x)
    if visible(y) then
      if not run then
        if i > 0 then
          run, startTip = { crossing(f, x, px) }, true
        else
          run, startTip = {}, leftTip
        end
      end
      run[#run + 1] = { x, y }
    elseif run then
      run[#run + 1] = crossing(f, px, x)
      emit(style, run, startTip, true)
      run = nil
    end
    px = x
  end
  if run then emit(style, run, startTip, rightTip) end
end

local function close(u, v)
  return math.abs(u - v) < 1e-4 * (plot.hi - plot.lo)
end

-- Lapos lista a töréspontokból (számok és táblázatok vegyesen).
local function flatten(list, out)
  out = out or {}
  for _, v in ipairs(list) do
    if type(v) == 'table' then flatten(v, out) else out[#out + 1] = v end
  end
  return out
end

-- a, a + h, ..., b
function steps(a, b, h)
  local out = {}
  for k = 0, math.floor((b - a) / h + 0.5) do out[#out + 1] = a + k * h end
  return out
end

-- Az f függvény grafikonja [a, b]-n. A breaks töréspontokban (szakadások,
-- pólusok) a görbe megszakad; ott a függvényérték pont, a tőle eltérő
-- féloldali határértékek karikák. Az ends két karaktere a két végpont:
-- < és > a görbe továbbfut, [ ] végpont, ( ) kimaradó végpont, . semmi.
function plot.graph(name, a, b, breaks, ends, style)
  local f = plot.fn[name]
  local xs = { a, b }
  for _, x in ipairs(flatten(breaks)) do
    if x > a + EPS and x < b - EPS then xs[#xs + 1] = x end
  end
  table.sort(xs)
  for i = 1, #xs - 1 do
    curve(f, xs[i], xs[i + 1], style,
      i == 1 and ends:sub(1, 1) == '<', i == #xs - 1 and ends:sub(2, 2) == '>')
  end
  local marks = {}
  local function mark(kind, x, y)
    if visible(y) then marks[#marks + 1] = { kind, x, y } end
  end
  for i = 2, #xs - 1 do
    local x = xs[i]
    local l, r, v = f(x - EPS), f(x + EPS), f(x)
    local defined = visible(v)
    if not (defined and close(l, v)) then mark('hole', x, l) end
    if not (defined and close(r, v)) and not close(l, r) then mark('hole', x, r) end
    if defined and not (close(l, v) and close(r, v)) then mark('end dot', x, v) end
  end
  for _, e in ipairs({ { ends:sub(1, 1), a, 1 }, { ends:sub(2, 2), b, -1 } }) do
    local c, x, side = e[1], e[2], e[3]
    if c == '[' or c == ']' then mark('end dot', x, f(x)) end
    if c == '(' or c == ')' then mark('hole', x, f(x + side * EPS)) end
  end
  for _, m in ipairs(marks) do
    tex.sprint(string.format('\\addplot[%s, %s] coordinates {(%.5f,%.5f)};',
      m[1], style, m[2], m[3]))
  end
end

local sqrt = math.sqrt
plot.fn.f1 = function(x) return (2 * x - 1) / (5 - x) end
plot.fn.f2 = function(x) return (4 - 6 * x) / (7 + 2 * x) end
plot.fn.f3 = function(x)
  return x * (sqrt(x * x + 1) - sqrt(math.max(0, x * x - 3)))
end
plot.fn.f4 = function(x)
  return sqrt(x * x + x + 1) - sqrt(math.max(0, x * x - x - 3))
end
\end{luacode*}

% \setwindow{xmin}{xmax}{ymin}{ymax}: a következő ábra ablaka.
\newcommand{\setwindow}[4]{%
	\def\wxmin{#1}\def\wxmax{#2}\def\wymin{#3}\def\wymax{#4}%
	\directlua{plot.window(#1, #2, #3, #4)}}
% \fngraph[stílus]{függvény}{a}{b}{töréspontok}{végpontok}
\newcommand{\fngraph}[6][]{\directlua{plot.graph("#2", #3, #4,
		load("return {\luaescapestring{\detokenize{#5}}}")(), "#6",
		"\luaescapestring{\detokenize{#1}}")}}
% Függőleges és vízszintes aszimptoták; \hasymptote{y}{felirat}{hol}, ahol
% a hol a felirat helye: a bal vagy a jobb szélen, a vonal alatt vagy fölött.
\newcommand{\vasymptote}[1]{\draw[mirror] (#1, \wymin) -- (#1, \wymax);}
\newcommand{\hasymptote}[3]{\draw[mirror] (\wxmin, #1) -- (\wxmax, #1)
	node[asymptote label/#3, font=\figtickfont, inner sep=2pt,
		text=gray!70!black] {#2};}
\tikzset{
	asymptote label/left above/.style={pos=0.02, above right},
	asymptote label/left below/.style={pos=0.02, below right},
	asymptote label/right above/.style={pos=0.98, above left},
	asymptote label/right below/.style={pos=0.98, below left},
}

\title{Az 5. gyakorlat függvényeinek ábrái}
\subtitle{Analízis A1 -- segédanyag}
\author{Sándor Tibor}
\date{}

\begin{document}
\maketitle

\section{Racionális törtfüggvények}

\begin{figure}[ht]
	\centering
	\setwindow{-25}{35}{-12}{8}
	\begin{tikzpicture}
		\begin{axis}[fn, xtick={-20, -10, 5, 10, 20, 30},
				ytick={-10, -5, 5}]
			\hasymptote{-2}{$y = -2$}{left below}
			\vasymptote{5}
			\fngraph[fcolor]{f1}{-25}{35}{5}{<>}
		\end{axis}
	\end{tikzpicture}
	\caption{Az $f(x) = \dfrac{2x - 1}{5 - x}$ függvény grafikonja;
		$\lim\limits_{x \to \pm\infty} f(x) = -2$.}
	\label{fig:f1-infinity}
\end{figure}

\begin{figure}[ht]
	\centering
	\setwindow{2}{4.8}{0}{12}
	\begin{tikzpicture}
		\begin{axis}[zoom, xtick={2, 3, 4}, ytick={4, 7, 10}]
			\draw[mirror] (4, \wymin) -- (4, 7) -- (\wxmin, 7);
			\fngraph[fcolor]{f1}{2}{4.8}{}{<>}
			\addplot[end dot, fcolor] coordinates {(4, 7)};
		\end{axis}
	\end{tikzpicture}
	\caption{Az $f(x) = \dfrac{2x - 1}{5 - x}$ függvény grafikonja a $4$ körül;
		$\lim\limits_{x \to 4} f(x) = 7$.}
	\label{fig:f1-point}
\end{figure}

\begin{figure}[ht]
	\centering
	\setwindow{-30}{25}{-13}{7}
	\begin{tikzpicture}
		\begin{axis}[fn, xtick={-20, -10, 10, 20}, ytick={-10, -5, 5},
				extra x ticks={-3.5}, extra x tick labels={$-\frac{7}{2}$}]
			\hasymptote{-3}{$y = -3$}{left above}
			\vasymptote{-3.5}
			\fngraph[gcolor]{f2}{-30}{25}{-3.5}{<>}
		\end{axis}
	\end{tikzpicture}
	\caption{A $g(x) = \dfrac{4 - 6x}{7 + 2x}$ függvény grafikonja;
		$\lim\limits_{x \to \pm\infty} g(x) = -3$.}
	\label{fig:f2-infinity}
\end{figure}

\begin{figure}[ht]
	\centering
	\setwindow{-2.6}{1.4}{-1}{8}
	\begin{tikzpicture}
		\begin{axis}[zoom, axis lines=middle, xtick={-2, -1, 1}, ytick={2, 4, 6}]
			\draw[mirror] (-1, 0) -- (-1, 2) -- (0, 2);
			\fngraph[gcolor]{f2}{-2.6}{1.4}{}{<>}
			\addplot[end dot, gcolor] coordinates {(-1, 2)};
		\end{axis}
	\end{tikzpicture}
	\caption{A $g(x) = \dfrac{4 - 6x}{7 + 2x}$ függvény grafikonja a $-1$
		körül; $\lim\limits_{x \to -1} g(x) = 2$.}
	\label{fig:f2-point}
\end{figure}

\section{Gyökös kifejezések}

\begin{figure}[ht]
	\centering
	\setwindow{-12}{12}{-4.5}{4.5}
	\begin{tikzpicture}
		\begin{axis}[fn, xtick={-10, -5, 5, 10}, ytick={-4, -2, 2, 4},
				extra x ticks={-1.7321, 1.7321},
				extra x tick labels={$-\sqrt{3}$, $\sqrt{3}$},
				extra x tick style={ticklabel pos=upper}]
			\hasymptote{2}{$y = 2$}{left above}
			\hasymptote{-2}{$y = -2$}{right below}
			\draw[mirror] (-1.7321, 0) -- (-1.7321, -3.4641);
			\draw[mirror] (1.7321, 0) -- (1.7321, 3.4641);
			\fngraph[fcolor]{f3}{-12}{-math.sqrt(3)}{}{<]}
			\fngraph[fcolor]{f3}{math.sqrt(3)}{12}{}{[>}
		\end{axis}
	\end{tikzpicture}
	\caption{Az $f(x) = x \left( \sqrt{x^2 + 1} - \sqrt{x^2 - 3} \right)$
		függvény grafikonja ($|x| \ge \sqrt{3}$);
		$\lim\limits_{x \to -\infty} f(x) = -2$ és
		$\lim\limits_{x \to \infty} f(x) = 2$.}
	\label{fig:f3}
\end{figure}

\begin{figure}[ht]
	\centering
	\setwindow{-12}{12}{-1.8}{3.6}
	\begin{tikzpicture}
		\begin{axis}[fn, xtick={-10, -5, 5, 10}, ytick={-1, 1, 2, 3}]
			\hasymptote{1}{$y = 1$}{left above}
			\hasymptote{-1}{$y = -1$}{right below}
			\fngraph[gcolor]{f4}{-12}{(1 - math.sqrt(13)) / 2}{}{<]}
			\fngraph[gcolor]{f4}{(1 + math.sqrt(13)) / 2}{12}{}{[>}
		\end{axis}
	\end{tikzpicture}
	\caption{A $g(x) = \sqrt{x^2 + x + 1} - \sqrt{x^2 - x - 3}$ függvény
		grafikonja ($x \le \frac{1 - \sqrt{13}}{2}$ vagy
		$x \ge \frac{1 + \sqrt{13}}{2}$);
		$\lim\limits_{x \to -\infty} g(x) = -1$ és
		$\lim\limits_{x \to \infty} g(x) = 1$.}
	\label{fig:f4}
\end{figure}

\end{document}
